\section{Related Work}
\label{sec:related}

\begin{table*}[t]
\centering
\caption{Comparison with evidence-aware multimodal verification and retrieval studies. An output of a set of different sizes of evidence is called adaptive; reranking, filtering, and modality routing are not adaptive set decoding.}
\label{tab:related-comparison}
\small
\setlength{\tabcolsep}{4pt}
\renewcommand{\arraystretch}{1.13}
\begin{tabular}{@{}>{\raggedright\arraybackslash}p{0.12\textwidth}
>{\raggedright\arraybackslash}p{0.16\textwidth}
>{\raggedright\arraybackslash}p{0.28\textwidth}
>{\raggedright\arraybackslash}p{0.14\textwidth}
>{\raggedright\arraybackslash}p{0.17\textwidth}@{}}
\toprule
Study & Setting & Retrieval or selection mechanism & Set policy & Evaluated outputs \\
\midrule
MOCHEG \cite{yao2023mocheg} & Web fact-checking & Separate neural text and image retrieval & Fixed top-$k$ & Retrieval, verdict, explanation \\
RED-DOT \cite{papadopoulos2025reddot} & VERITE; NewsCLIPpings+ & Relevance detection and evidence reranking & Fixed pool with item masking & Item relevance, verdict \\
SciVer \cite{wang2025sciver} & Scientific papers & Foundation-model and RAG benchmark & No adaptive set policy & Evidence, verdict \\
MuSciClaims \cite{lal2025musciclaims} & Scientific figures & Figure localization and reasoning diagnostics & Provided figure context & Localization, verdict \\
M$^2$-TabFact \cite{zhou2025m2tabfact} & Multi-document tables & Textual and visual table evidence reasoning & Provided multimodal evidence & Verdict \\
MEVER \cite{zhang2026mever} & MOCHEG; FACTIFY; AIChartClaim & Graph retrieval and multimodal fusion & Fixed top-$k$ & Retrieval, verdict, explanation \\
Utility-oriented \cite{luo2026utility} & MRAG-Bench; Visual-RAG & Information-gain and helpfulness ranking & Fixed top-$K$ & Downstream answer utility \\
Selective retrieval \cite{tsoneva2026selective} & AVerImaTeC & Question-type routing with hybrid retrieval & Fixed per-question budget & Evidence, verdict, justification \\
\midrule
\textbf{This study} & \textbf{MOCHEG and AVerImaTeC organizer stores} & \textbf{Label-blind multimodal packets, calibrated membership scores, packet-completeness estimate} & \textbf{Fixed, threshold, expected-F1, and switch controls} & \textbf{Reachability, set F1, calibration, shared verifier, cross-benchmark transfer} \\
\bottomrule
\end{tabular}
\end{table*}

\subsection{Evidence-Based Fact-Checking and Evaluation}

Evidence-based verification can typically be broken down into four components: claim analysis, evidence retrieval, evidence assessment and verdict prediction \cite{zeng2021automated,zhou2020survey,islam2020survey}. The findings vary with the domain, label construction, modality coverage, and controls for event or source overlap, as documented in dataset surveys \cite{dulizia2021datasets,rastogi2022threet}. Difficulties with cross-domain transfer, deficient explanations, and benchmark protocols remain common \cite{mridha2021review,aimeur2023review,lv2025survey}. Our assessment therefore follows the retrieval-and-selection chain rather than content-only classification.

Evidence requirements are revealed through text benchmarks. SciFact contains sentence evidence from scientific publications, HoVer is multi-document reasoning, COVID-Fact is following evolving claims in the public-health arena, and X-Fact introduces multilingual variation \cite{wadden2020scifact,jiang2020hover,saakyan2021covidfact,gupta2021xfact}. External-document systems minimize reliance on stylistic shortcuts, while also incurring retrieval, freshness, and attribution errors \cite{zhang2023realtime}. So, we separate evidence-set quality and verdict accuracy.

\subsection{Multimodal Misinformation Detection}

Content-based multimodal detectors include bilinear pooling \cite{kumari2021amfb}, hierarchical attention \cite{meel2021han}, topic memory \cite{ying2021topicmemory} and cross-attention \cite{ying2021crossattention}, ambiguity learning \cite{chen2022cafe} and multi-level fusion \cite{jing2023mpfn} detectors. Contrastive, in-modality, stylometric or semantic-alignment signals are added in later systems \cite{yan2024fmc,zhu2024ifis,nadeem2023ssm,uppada2022imagetext}. These designs help to classify at the post level; however, the fused representation does not reveal the source of a verdict, such as an external document or image.

Recent works focus on local-global mismatch \cite{wu2023mfir}, cross-modal knowledge distillation \cite{hu2024matching,peng2024contextual}, contextual inconsistency \cite{shang2026semantic} and conflicting unimodal signals \cite{wang2024neutralization}. Out-of-context benchmarks like NewsCLIPpings \cite{luo2021newsclippings}, COSMOS \cite{aneja2023cosmos}, and VERITE \cite{papadopoulos2024verite} evaluate pairs of images and captions where they do not match in context but control unimodal bias. DGM4, SNIFFER, MMDFND, and IMFND extend localization, explanation, domain coverage or vision-language model usage \cite{shao2023dgm4,qi2024sniffer,tong2024mmdfnd,jiang2025imfnd}. They are not primarily aimed at a set of external evidence identifiers that can be audited, but at manipulation or veracity.

\subsection{Multimodal Evidence Retrieval and Verification}

MOCHEG offers a common evaluation criterion for multimodal retrieval, three-way verdict prediction and explanation generation \cite{yao2023mocheg}. MEVER builds a text--image evidence graph and performs token- and evidence-level fusion \cite{zhang2026mever}, while RED-DOT identifies and reranks external multimodal evidence before prediction \cite{papadopoulos2025reddot}. Other recent studies distinguish utility from semantic similarity and learn when visual retrieval is needed \cite{luo2026utility,tsoneva2026selective}. These systems address fixed retrieval, relevance filtering, modality routing, utility ranking, and adaptive set decoding as different decision problems.

AVerImaTeC assesses real-world image--text claims using open-web or organizer-provided knowledge, and its shared task conditions verdict scoring on evidence quality \cite{cao2026averimatecshared}. In that context, REVEAL combines iterative sparse and dense retrieval \cite{tariq2026reveal}. MOCHEG's stable item corpus differs from question--answer evidence and open-web access, so transfer requires a separately frozen unit definition and candidate-construction protocol.

\subsection{Scientific and Citation-Grounded Evidence}

Stable evidence locations in passages, tables and figures are provided by scientific verification. SciVer consists of claim-evidence instances extracted from scientific papers and annotated by experts; MuSciClaims focuses on figure-grounded claims and controlled reasoning perturbations; and M$^2$-TabFact covers tables presented across text and images \cite{wang2025sciver,lal2025musciclaims,zhou2025m2tabfact}. These tools can be used to make more precise localization analysis, but they do not explicitly reflect political or breaking-news verification.

An explanation might be expressed effectively, but with the use of the incorrect source. Evidence-rationale benchmarks check whether a prediction is connected to a selected span, and faithful-by-construction rationalizers limit predictors to explicit evidence subsets \cite{deyoung2020eraser,jain2020faithful}. Citation-grounded generation assesses whether retrieved sources support generated statements \cite{gao2023alce}; RAGTruth catalogs hallucinations that persist despite retrieval \cite{niu2024ragtruth}; and citation masking tests whether a fact-checking explanation recovers the correct evidence \cite{xing2025attribution}. The present study adopts released evidence identifiers rather than generated rationales and allows exact-set precision, recall, F1, Jaccard similarity and exact match.

\subsection{Calibration and Adaptive Set Decisions}

A score may rank a candidate correctly while remaining a poor probability. Post-hoc calibration includes temperature and Platt mappings, beta calibration, and nonparametric isotonic regression; each makes different assumptions about the distortion being corrected \cite{silvafilho2023calibration}. After class weighting, prior correction is important because the optimized training distribution differs from the operating distribution. The calibrated membership distribution can then support expected-utility decoding, while structurally different policies such as top-$k$, thresholding, and learned count use the same scores.

In this study, representation models and evidence identities remain fixed while a bounded, label-blind rank fusion changes only the retained ordering. The selector experiment then evaluates packet reachability, prior-aware probability correction, exact expected-F1 decoding, and one prespecified completeness-aware switch under matched candidates, tuning roles, and downstream verification. A separately frozen organizer-store run tests the same relationships on AVerImaTeC without transferring fitted MOCHEG parameters. This design isolates which conclusions are benchmark-specific and which survive a change in evidence units, candidate construction, and class balance (Table~\ref{tab:related-comparison}).
